% Part: first-order-logic
% Chapter: tableaux
% Section: rules-and-proofs

\documentclass[../../../include/open-logic-section]{subfiles}

\begin{document}

\iftag{FOL}
      {\olfileid{fol}{tab}{rul}}
      {\olfileid{pl}{tab}{rul}}

\olsection{Regels en \usetoken{P}{tableau}}

!!^a{tableau} zet systematisch de mogelijke manieren op een rij
waarop !!a{sentence} in !!a{structure} waar of onwaar kan zijn. Een
tableau bestaat uit !!{signed formula}s. Dat zijn !!{sentence}s met
een waarheidswaarde als ``teken'': $\True$ of~$\False$. Deze
!!{formula}s staan in een boom die naar beneden groeit.

\begin{defn}
  Een \emph{!!{signed formula}} is een paar. Het paar bestaat uit een
  waarde, waar of onwaar, en !!a{sentence}. Er zijn dus twee
  mogelijkheden:
  \[
  \sFmla{\True}{!A} \text{ of } \sFmla{\False}{!A}.
  \]
\end{defn}

We kunnen $\sFmla{\True}{!A}$ lezen als ``$!A$ zou waar kunnen zijn''
en $\sFmla{\False}{!A}$ als ``$!A$ zou onwaar kunnen zijn'' (in een
zekere !!{structure}).

Elke !!{signed formula} in de boom heeft een van twee rollen. De
formule is een \emph{aanname}; de aannamen staan helemaal bovenaan in
de boom. Of de formule volgt met een afleidingsregel uit !!a{signed
formula} die hoger in de boom staat. Voor elke mogelijke
!!{main operator} van die eerdere
!!{formula} bestaan twee regels. De ene geldt als het teken~$\True$
is, de andere als het teken~$\False$ is. Door sommige regels splitst
de boom zich. Andere regels voegen alleen !!{signed formula}s aan de
tak toe. Je past een regel niet per se toe op de !!{signed formula}
die er direct boven staat. Je mag de regel toepassen op een
willekeurige !!{signed formula} die op de tak staat tussen de wortel
en die plaats; vaak moet dat zelfs.

Een tak is \emph{gesloten} als hij zowel $\sFmla{\True}{!A}$ als
$\sFmla{\False}{!A}$ bevat. Een !!{tableau} is gesloten als elke tak
gesloten is. Elke tak beschrijft één manier waarop de formules erop
samen zouden kunnen kloppen. Maar $\sFmla{\True}{!A}$ en
$\sFmla{\False}{!A}$ kunnen niet tegelijk kloppen. Een gesloten tak
beschrijft dus een mogelijkheid die niet kan bestaan. Een gesloten
!!{tableau} sluit daarom elke mogelijkheid uit
waarin tegelijk alle aannamen van de vorm $\sFmla{\True}{!A}$ waar en
alle aannamen van de vorm~$\sFmla{\False}{!A}$ onwaar zijn.

Een gesloten !!{tableau} \emph{voor $!A$} is een gesloten
!!{tableau} met $\sFmla{\False}{!A}$ als wortel. Als zo'n gesloten
!!{tableau} bestaat, kan $!A$ op geen enkele manier onwaar zijn. Dus
moet $!A$ waar zijn in elke !!{structure}.

\end{document}
